\clearpage
\section{The Two Great Pentecostal Hymns}
\label{sec:hymns}

These are complete received Latin texts as printed in the \textit{Compendium of the Catechism of the Catholic Church}. The facing English texts are original literary translations for this project, not official liturgical translations. In private devotion one may pray either language or both. In a liturgical celebration use only the text and language authorized for that celebration.

\subsection{\textit{Veni Creator Spiritus}}

The guide appoints this hymn every day because its movement is a compact theology of invocation: Creator and Gift; light and love; strength and peace; knowledge of Father and Son; Trinitarian praise. The traditional Rabanus Maurus attribution is possible but uncertain; the hymn belongs to the ninth-century Carolingian world.

\bilingualprayer{Come, Creator Spirit}{Received Latin hymn; project literary translation.}{%
Veni, creátor Spíritus,\\
mentes tuórum vísita,\\
imple supérna grátia\\
quæ tu creásti péctora.\\[0.55em]
Qui díceris Paráclitus,\\
altíssimi donum Dei,\\
fons vivus, ignis, cáritas,\\
et spiritális únctio.\\[0.55em]
Tu septifórmis múnere,\\
dígitus patérnæ déxteræ,\\
tu rite promíssum Patris,\\
sermóne ditans gúttura.\\[0.55em]
Accénde lumen sénsibus,\\
infúnde amórem córdibus,\\
infírma nostri córporis\\
virtúte firmans pérpeti.\\[0.55em]
Hostem repéllas lóngius\\
pacémque dones prótinus;\\
ductóre sic te prævio\\
vitémus omne nóxium.\\[0.55em]
Per te sciámus da Patrem,\\
noscámus atque Fílium,\\
te utriúsque Spíritum\\
credámus omni témpore.\\[0.55em]
Deo Patri sit glória,\\
et Fílio, qui a mórtuis\\
surréxit, ac Paráclito,\\
in sæculórum sæcula. Amen.}{%
Come, Spirit, Maker; visit now\\
the minds and hearts that are thine own;\\
fill with the grace that comes from heaven\\
the breasts whose life was made by thee.\\[0.55em]
Thou who art named the Paraclete,\\
Gift of the God who is Most High,\\
living spring, fire, charity,\\
and holy anointing of the soul.\\[0.55em]
Sevenfold in the grace thou givest,\\
finger of the Father's right hand,\\
promise pledged to us by the Father,\\
enrich our tongues with living speech.\\[0.55em]
Kindle thy light within our senses;\\
pour thy love into our hearts;\\
strengthen the frailty of our bodies\\
with power that can never fail.\\[0.55em]
Drive far away the ancient foe,\\
and grant us peace without delay;\\
with thee before us as our guide,\\
may we escape all hurtful things.\\[0.55em]
Through thee grant us to know the Father,\\
and learn the Son whom he has sent;\\
and thee, the Spirit of them both,\\
confess in faith through every age.\\[0.55em]
To God the Father glory be,\\
and to the Son, from death arisen,\\
and to the Paraclete, all praise\\
through ages unto ages. Amen.}

\subsection{\textit{Veni Sancte Spiritus}}

This sequence belongs to the Roman Mass of Pentecost. Its authorship is uncertain; early-thirteenth-century witnesses and competing attributions counsel restraint. The hymn asks the Spirit to act within created poverty. Its verbs form an examination: wash, water, heal, bend, warm, and guide. Day~9 uses it in full as the threshold of the feast.

\bilingualprayer{Come, Holy Spirit}{Received Latin Pentecost sequence; project literary translation.}{%
Veni, Sancte Spíritus,\\
et emítte cælitus\\
lucis tuæ rádium.\\[0.45em]
Veni, pater páuperum,\\
veni, dator múnerum,\\
veni, lumen córdium.\\[0.45em]
Consolátor óptime,\\
dulcis hospes ánimæ,\\
dulce refrigérium.\\[0.45em]
In labóre réquies,\\
in æstu tempéries,\\
in fletu solácium.\\[0.45em]
O lux beatíssima,\\
reple cordis íntima\\
tuórum fidélium.\\[0.45em]
Sine tuo númine,\\
nihil est in hómine,\\
nihil est innóxium.\\[0.45em]
Lava quod est sórdidum,\\
riga quod est áridum,\\
sana quod est sáucium.\\[0.45em]
Flecte quod est rígidum,\\
fove quod est frígidum,\\
rege quod est dévium.\\[0.45em]
Da tuis fidélibus,\\
in te confidéntibus,\\
sacrum septenárium.\\[0.45em]
Da virtútis méritum,\\
da salútis éxitum,\\
da perénne gáudium. Amen.}{%
Come, Holy Spirit, and send\\
from the high court of heaven\\
one radiant beam of thy light.\\[0.45em]
Come, Father of all the poor;\\
come, generous Giver of gifts;\\
come, hidden Light of every heart.\\[0.45em]
Thou best of all consolers,\\
the soul's most gentle guest,\\
her sweetness and refreshment.\\[0.45em]
In toil, be thou our rest;\\
in fevered heat, our coolness;\\
amid our tears, our solace.\\[0.45em]
O light of blessedness,\\
fill the inmost chambers\\
of the hearts of thy faithful.\\[0.45em]
Without thy sovereign presence\\
there is nothing whole in man,\\
nothing that remains untainted.\\[0.45em]
Wash whatever has been soiled;\\
water whatever has gone dry;\\
heal whatever has been wounded.\\[0.45em]
Bend whatever has grown rigid;\\
warm whatever has grown cold;\\
guide whatever has gone astray.\\[0.45em]
Give unto thy faithful,\\
who place their trust in thee,\\
the holy gift in sevenfold grace.\\[0.45em]
Give the ripened work of virtue;\\
give the journey's saving end;\\
give joy that has no ending. Amen.}

\begin{studybox}{Copyright and use note}
The ancient Latin texts are reproduced as governing prayer witnesses. The English versions were composed for this document and may be revised after specialist review. No modern proprietary Missal translation is reproduced. Parish or liturgical use requires the competent approved book, not these project translations.
\end{studybox}
